% Generated draft: requires mathematical and source review.
\documentclass[11pt]{article}
\usepackage[margin=25mm]{geometry}
\usepackage{fontspec}
\setmainfont{Latin Modern Roman}
\usepackage{amsmath,amssymb,mathtools,unicode-math}
\setmathfont{Latin Modern Math}
\usepackage{enumitem,needspace,xurl,hyperref}
\setlength{\parindent}{0pt}
\setlength{\parskip}{5pt}
\setlength{\emergencystretch}{3em}
\newcommand{\N}{\mathbb{N}}
\newcommand{\Z}{\mathbb{Z}}
\newcommand{\Q}{\mathbb{Q}}
\newcommand{\R}{\mathbb{R}}
\newcommand{\C}{\mathbb{C}}
\newcommand{\F}{\mathbb{F}}
\newcommand{\A}{\mathbb{A}}
\newcommand{\PP}{\mathbb{P}}
\newcommand{\E}{\mathbb{E}}
\newcommand{\eps}{\varepsilon}
\newcommand{\dd}{\,\mathrm{d}}
\begin{document}
\section*{UnsolvedMath ID: 701180}
\textbf{Source problem code:} catalogue entry
\subsection{Definitions and mathematical statement}
For positive integers \(l,m,n\), retain the presentation exactly as supplied:
\[
T_{l,m,n}=\langle x,y\mid x^l=y^m=(xy)^n\rangle
=\langle x,y\mid x^l y^{-m}=1,\ y^m(xy)^{-n}=1\rangle.
\]
The displayed relations equate three powers; they do not require their common value to be the identity. In particular, this presentation must be distinguished from
\[
\Delta_{l,m,n}=\langle x,y\mid x^l=y^m=(xy)^n=1\rangle.
\]
The source calls its displayed group a triangle group, but does not clarify whether the missing relation setting the powers equal to the identity is intentional. The literal presentation \(T_{l,m,n}\) is retained below; replacing it by \(\Delta_{l,m,n}\) gives a separate possible intended formulation.

For a specified class \(\mathcal A\) of finite alternating groups, say that a group \(H\) is residually \(\mathcal A\) through epimorphisms if
\[
\forall h\in H\setminus\{1\}\ \exists A\in\mathcal A\ \exists\text{ a surjective homomorphism }\phi:H\longrightarrow A
\quad\text{such that}\quad \phi(h)\ne1.
\]
Here \(A_d\) denotes the group of even permutations of \(\{1,\ldots,d\}\). Two possible choices are
\[
\mathcal A_{\mathrm{all}}=\{A_d:d\in\Z,\ d\ge1\},
\qquad
\mathcal A_{\mathrm{simple}}=\{A_d:d\in\Z,\ d\ge5\}.
\]
The source does not specify the allowed degrees or whether the separating homomorphisms must be surjective. Thus this definition makes an explicit conditional interpretation of the source phrase ``residually finite alternating groups''; it does not resolve that terminology.

With these qualifications, the mathematical target is the following question, for the intended choice of \(\mathcal A\): is it true that for every integer \(n>6\) there exist an uncountable set \(I\) and normal subgroups \(N_i\trianglelefteq T_{2,3,n}\), indexed by \(i\in I\), such that
\[
T_{2,3,n}/N_i\text{ is residually }\mathcal A
\quad\text{for every }i\in I,
\]
with residuality understood through epimorphisms as above, and
\[
T_{2,3,n}/N_i\not\cong T_{2,3,n}/N_j
\qquad\text{whenever }i\ne j?
\]
The source reports that the corresponding assertion for parameters \((2,3,30)\) was proved by A. M. Brunner, R. G. Burns and J. Wiegold in 1979. That reported result inherits the presentation and terminology qualifications above.

The supplied registry labels this problem open. Its background also records an explicitly unverified affirmative AI solution claim attributed to M. Kinyon and J. Urban in a 2026 preprint. Neither the label nor that claim is an independently verified mathematical conclusion here.
\subsection{Short English statement}
J. Wiegold asks whether, for every integer \(n>6\), the group with parameters \((2,3,n)\) has uncountably many pairwise nonisomorphic quotients that are residually finite alternating. The source reports the case \(n=30\). Its presentation omits relations setting the displayed powers equal to the identity, and its residuality terminology needs clarification. The supplied open label coexists with an explicitly unverified 2026 affirmative AI solution claim.
\subsection{Sources}
\begin{itemize}
\item Local UnsolvedMath registry, supplied complete source record: permanent ID \(701180\), exact source problem code \texttt{catalogue entry}, title \emph{Kourovka Notebook Problem 15.92}. Proposed by J. Wiegold in 2002. The supplied status is \texttt{open}.
\item \emph{The Kourovka Notebook (Unsolved Problems in Group Theory)}, 15th issue (2002), Problem 15.92; also identified in the record with the 21st issue edition (2026), arXiv:1401.0300v48. The supplied collection description names E. I. Khukhro and V. D. Mazurov as editors.
\item A. M. Brunner, R. G. Burns and J. Wiegold, \emph{Math. Scientist} \textbf{4} (1979), 93--98. The supplied statement attributes the result for parameters \((2,3,30)\) to this reference; no article title is supplied.
\item M. Kinyon and J. Urban, \emph{Preprint}, 2026, cited in the supplied background as an ``unverified AI solution'' claiming an affirmative answer: \url{https://kourovkanotebookorg.wordpress.com/wp-content/uploads/2026/09/kourovka-mathematics.pdf}. This attribution and claim have not been independently verified.
\item The supplied status review is dated \texttt{2026-10-16}, records canonical status \texttt{open} and evidence \texttt{unrefereed\_claim}, and states that the problem is marked with \(\#\) for the unverified AI claim in the October 2026 21st issue. This assessment is supplied dated curation, not independently verified.
\end{itemize}
\end{document}
